
We give a small calculus for a single agent whose only source of nondeterminism is calls to
a language model, together with a journal that records every outcome.
The calculus is deliberately minimal: it isolates the part of an agent that durable
execution has to get right, namely the interplay between stochastic model calls and the
record that makes them reproducible.

\subsection{Syntax}
\label{sec:syntax}

Let $\mathit{Val}$ be a countable set of values; prompts and model outputs are values (for
example, finite token sequences).
Expressions $e$ are pure and deterministic: $e \Downarrow v$ is a total function from closed
expressions to values, and we assume programs are well typed, so the condition of an
$\mathbf{if}$ evaluates to a boolean.
Terms are
\[
  M, N \;::=\; \mathbf{return}\;e
    \;\mid\; \mathbf{let}\;x = e\;\mathbf{in}\;M
    \;\mid\; \mathbf{if}\;e\;\mathbf{then}\;M\;\mathbf{else}\;N
    \;\mid\; \mathbf{let}\;x \Leftarrow \mathbf{oracle}(e)\;\mathbf{in}\;M .
\]
The term $\mathbf{let}\;x \Leftarrow \mathbf{oracle}(e)\;\mathbf{in}\;M$ asks the model with
the prompt that $e$ evaluates to and binds the response to $x$ in $M$.
An \emph{oracle} assigns to each prompt $p$ a probability measure $\meas{p}$ on
$\mathit{Val}$; since $\mathit{Val}$ is countable, $\meas{p}$ is discrete.
Nothing else is assumed about $\meas{p}$: the results below hold for every oracle.
The same construct models any externally determined value, such as a tool response, a
message from the environment, or a clock read; we speak of model calls for concreteness.
We omit general recursion, so every run makes finitely many oracle calls
(Section~\ref{sec:limitations} discusses lifting this).

\paragraph{Programs as trees.}
Because expressions are pure, the behaviour of a closed term is captured by its tree of
oracle calls: either it halts with a value, or it makes a call with some prompt $p$ and
continues with a subtree that depends on the response.
Formally, define $T(M)$ by
\begin{align*}
  T(\mathbf{return}\;e) &= \mathsf{halt}(v) && \text{where } e \Downarrow v,\\
  T(\mathbf{let}\;x = e\;\mathbf{in}\;M) &= T(M[v/x]) && \text{where } e \Downarrow v,\\
  T(\mathbf{if}\;e\;\mathbf{then}\;M\;\mathbf{else}\;N) &= T(M) \text{ or } T(N) && \text{as } e \Downarrow \mathsf{true} \text{ or } \mathsf{false},\\
  T(\mathbf{let}\;x \Leftarrow \mathbf{oracle}(e)\;\mathbf{in}\;M) &= \mathsf{draw}\bigl(p,\; v \mapsto T(M[v/x])\bigr) && \text{where } e \Downarrow p.
\end{align*}
Trees are well founded because terms have no recursion.
All definitions and results below depend only on $T(M)$, and the Lean development
(Section~\ref{sec:mechanization}) works directly with these trees.

\subsection{The Journal and Three Modes}
\label{sec:modes}

A \emph{journal} $J$ is a finite list of entries $(c, p, v)$: the cursor $c$ (the entry's
position in the log), the prompt $p$, and the response $v$.
The agent runs in one of three modes.
\begin{description}[noitemsep]
  \item[Record.] Each call is answered by the oracle, and its entry is appended to the journal.
  \item[Replay.] Each call is answered from the next journal entry; the oracle is not consulted.
  \item[Resume.] Replay until the journal is exhausted, then switch to record mode for the rest
    of the run. This is recovery after a crash.
\end{description}
A run emits a \emph{trace}: one label $o(c, p, v)$ per call, recording its cursor, prompt,
and response, followed by $\mathsf{done}(v)$ when the run halts with $v$.

\paragraph{Randomness tapes.}
To compare runs in different modes we fix their randomness in advance.
A \emph{tape} is a function $\omega : \mathbb{N} \times \mathit{Val} \to \mathit{Val}$, where
$\omega(c, p)$ is the response the oracle gives to prompt $p$ at cursor $c$.
Let $\Omega$ be the set of tapes with the product measure
$\mathbb{P} = \bigotimes_{(c,p)} \meas{p}$, under which the coordinates are independent and
$\omega(c,p)$ has law $\meas{p}$.
Since $\mathit{Val}$ is countable, so is the set of finite traces, so every set of traces is
measurable and every function of a run defined below is a random variable.

\paragraph{Semantics.}
Write $\mathsf{rec}(T, \omega, c)$, $\mathsf{rep}(T, J)$, and $\mathsf{res}(T, J, \omega, c)$
for the (trace, journal) pair produced by running tree $T$ from cursor $c$ in each mode:
\begin{align*}
  \mathsf{rec}(\mathsf{halt}(v), \omega, c) &= (\mathsf{done}(v),\; \varepsilon)\\
  \mathsf{rec}(\mathsf{draw}(p, k), \omega, c) &= o(c,p,v) \cdot \mathsf{rec}(k(v), \omega, c{+}1)
    && v = \omega(c,p), \text{ entry } (c,p,v)\\[3pt]
  \mathsf{rep}(\mathsf{halt}(v), J) &= (\mathsf{done}(v),\; \varepsilon)\\
  \mathsf{rep}(\mathsf{draw}(p, k), (c',p',v) \cdot J) &= o(c',p,v) \cdot \mathsf{rep}(k(v), J)
    && \text{entry } (c',p,v)\\
  \mathsf{rep}(\mathsf{draw}(p, k), \varepsilon) &= (\varepsilon, \varepsilon)
    && \text{journal exhausted}
\end{align*}
Here $\ell \cdot (t, J)$ prepends the label $\ell$ to the trace $t$ and the indicated entry to
$J$; replay re-emits each journal entry it consumes, so its output journal can be compared
with the recorded one.
Resume mode replays while the journal lasts and then records, tracking the cursor $c$ so that
the first fresh call lands at the position after the last surviving entry:
\begin{align*}
  \mathsf{res}(\mathsf{halt}(v), J, \omega, c) &= (\mathsf{done}(v),\; \varepsilon)\\
  \mathsf{res}(\mathsf{draw}(p, k), (c',p',v) \cdot J, \omega, c) &= o(c',p,v) \cdot \mathsf{res}(k(v), J, \omega, c{+}1)
    && \text{entry } (c',p,v)\\
  \mathsf{res}(\mathsf{draw}(p, k), \varepsilon, \omega, c) &= o(c,p,v) \cdot \mathsf{res}(k(v), \varepsilon, \omega, c{+}1)
    && v = \omega(c,p), \text{ entry } (c,p,v)
\end{align*}
These functions correspond to the small-step rules
\begin{align*}
  &\textsc{[O-rec]} && \langle \mathsf{draw}(p,k), J, c \rangle \xrightarrow{o(c,p,v)} \langle k(v), J\cdot(c,p,v), c{+}1 \rangle
     && \text{with probability } \meas{p}(\{v\}),\\
  &\textsc{[O-rep]} && \langle \mathsf{draw}(p,k), (c',p',v)\cdot R, c \rangle \xrightarrow{o(c',p,v)} \langle k(v), R, c{+}1 \rangle
     && \text{with probability } 1,\\
  &\textsc{[O-resume]} && \langle \mathsf{draw}(p,k), \varepsilon, c \rangle \text{ in resume mode steps as \textsc{[O-rec]}},
\end{align*}
where in replay and resume mode $R$ is the part of the journal not yet replayed.
In replay mode an exhausted journal at a call is stuck; in resume mode it falls through to
recording.
\textsc{[O-rec]} is the only rule whose successor is not a single configuration, so all
randomness in a run comes from oracle calls.
Replay checks neither the journaled prompt nor the cursor against the program: it trusts that
it is replaying the same program from the same starting cursor (assumption A2 below).

\begin{figure}[t]
  \centering
  \fitwidth{\begin{tikzpicture}[
  >=Stealth,
  every node/.style={align=center, font=\small},
  term/.style={draw, rounded corners=2pt, minimum width=4.8cm, minimum height=0.75cm, inner sep=3pt},
  state/.style={draw, rounded corners=2pt, minimum width=2.7cm, minimum height=0.62cm, inner sep=2pt, font=\footnotesize},
  jbox/.style={draw, densely dashed, rounded corners=1pt, minimum width=3.6cm, minimum height=0.58cm, inner sep=3pt, font=\footnotesize},
  lbl/.style={font=\scriptsize\itshape},
  arrow/.style={->, thick},
  darrow/.style={->, semithick, densely dashed},
]
\def\xR{-4.7}\def\xP{4.7}

\node[term] (draw) at (0,0) {$\texttt{let}\;x \leftarrow \mathrm{oracle}(e)\;\texttt{in}\;M$};
\node[font=\bfseries] at (\xR,1.2) {RECORD mode};
\node[font=\bfseries] at (\xP,1.2) {REPLAY mode};
\draw[very thick, dotted] (0,-0.7) -- (0,-5.5);

\coordinate (rb) at (\xR,-1.5);
\draw[arrow] (draw.south west) -- node[lbl, above left=-2pt]{$[O\text{-}rec]$} (rb);
\node[state] (v1) at (\xR-1.7,-3.0) {$C'_1\,[x\mapsto v_1]$};
\node[state] (v2) at (\xR+1.7,-3.0) {$C'_2\,[x\mapsto v_2]$};
\node at (\xR+3.3,-3.0) {$\cdots$};
\draw[arrow] (rb) -- node[lbl, left=1pt]{$\mu_p(\{v_1\})$} (v1.north);
\draw[arrow] (rb) -- node[lbl, right=1pt]{$\mu_p(\{v_2\})$} (v2.north);
\node[jbox] (jrec) at (\xR,-4.3) {$J \mathrel{+\!=} (c, p, v_i)$};
\draw[darrow] (v2.south) -- (jrec.north);
\node[lbl, text width=6.2cm] at (\xR,-5.15)
  {probabilistic successor distribution $\mathcal{D}_{rec}$; the realized $v_i$ is appended to $J$};

\coordinate (pb) at (\xP,-1.5);
\draw[arrow] (draw.south east) -- node[lbl, above right=-2pt]{$[O\text{-}rep]$} (pb);
\node[state] (vr) at (\xP,-3.0) {$C'\,[x\mapsto v_j]$};
\draw[arrow] (pb) -- node[lbl, right=1pt]{probability $=1$} (vr.north);
\node[jbox] (jrep) at (\xP,-4.3) {next entry $(c, p, v_j)$};
\draw[darrow] (vr.south) -- (jrep.north);
\node[lbl, text width=6.2cm] at (\xP,-5.15)
  {deterministic read (no append); same cursor $c$ as the record run (Lemma~\ref{lem:cursors})};
\end{tikzpicture}
}
  \caption{Record and replay of one oracle call.
    In record mode (left) the call samples $v \sim \meas{p}$ and appends $(c, p, v)$ to the
    journal. In replay mode (right) the same call reads the next entry with probability~1 and
    appends nothing new. The cursor $c$ does not depend on any sampled value
    (Lemma~\ref{lem:cursors}), so both modes address the same entry.}
  \label{fig:record-replay}
\end{figure}

\paragraph{The record distribution.}
The record run samples each response from the oracle as it goes.
Its distribution is that of $\mathsf{rec}(T, \omega, c)$ for $\omega \sim \mathbb{P}$:

\begin{lemma}[Tapes realize sampling]
\label{lem:tape}
Call $t = o(c,p_1,v_1) \cdots o(c{+}n{-}1,p_n,v_n)\,\mathsf{done}(v)$ a \emph{trace of $T$
from $c$} if it follows a path of $T$: each $p_i$ is the prompt $T$ issues after the responses
$v_1, \ldots, v_{i-1}$, and $T$ then halts with $v$.
For every tree $T$, cursor $c$, and finite trace $t$ of this form,
\[
  \mathbb{P}\bigl[\,\mathsf{rec}(T,\omega,c) \text{ has trace } t\,\bigr]
  \;=\; \prod_{i=1}^{n} \meas{p_i}(\{v_i\})
\]
if $t$ is a trace of $T$, and $0$ otherwise. This is the probability that the sampling
semantics (rule \textsc{[O-rec]}) produces $t$.
\end{lemma}

\begin{proof}
The $i$-th call of the run is made at cursor $c+i-1$ with a prompt $p_i$ that is determined by
the responses $v_1, \ldots, v_{i-1}$, and its response is the coordinate
$\omega(c+i-1, p_i)$.
The event that the run produces $t$ is therefore the event that
$\omega(c+i-1, p_i) = v_i$ for $i = 1, \ldots, n$, where $p_1, \ldots, p_n$ are the prompts
that $T$ issues along $t$ (fixed once $t$ is fixed).
These are $n$ distinct coordinates of the tape, so by independence the event has probability
$\prod_i \meas{p_i}(\{v_i\})$.
The sampling semantics produces $t$ with the same probability, by the chain rule:
given the first $i-1$ responses, the $i$-th is drawn from $\meas{p_i}$.
\end{proof}
